\documentclass[aspectratio=169,11pt]{beamer}
% Shared look for every deck: Metropolis with a progress bar under each frame title.
% Build with Tectonic: `make` in this directory.

\usetheme[progressbar=frametitle,numbering=fraction,sectionpage=progressbar,block=fill]{metropolis}

% Metropolis looks Fira Sans up by system name; load it from the TeX bundle instead,
% Light for body text as Metropolis intends, with lining figures for data.
\usepackage[sfdefault,light,lf]{FiraSans}
\usepackage{booktabs,tabularx}
\usepackage{tikz}
\usetikzlibrary{arrows.meta,positioning,shapes.geometric,calc}

% Metropolis's palette, with the decks' own colour names mapped onto it.
\definecolor{mDarkTeal}{HTML}{23373B}
\definecolor{mDarkBrown}{HTML}{604C38}
\definecolor{mLightBrown}{HTML}{EB811B}
\colorlet{crust}{mDarkBrown}
\definecolor{crumb}{HTML}{FAEEE2}
\colorlet{rye}{mDarkTeal}
\definecolor{flag}{HTML}{B03A2E}
\definecolor{own}{HTML}{2E7D4F}
\definecolor{muted}{HTML}{7A7F80}

% Top-align frames, so the \vfill in \source puts the source line at the bottom.
% A slimmer title bar and narrower margins than Metropolis's defaults leave room
% for the denser slides.
\makeatletter
\beamer@centeredfalse
\setlength{\metropolis@frametitle@padding}{1.4ex}
% A progress bar you can see: under every frame title and on section pages.
\setlength{\metropolis@progressinheadfoot@linewidth}{2pt}
\setlength{\metropolis@progressonsectionpage@linewidth}{1.5pt}
\setlength{\metropolis@titleseparator@linewidth}{1pt}
% Metropolis's title page is \paperheight tall, which overflows a top-aligned frame.
\setbeamertemplate{title page}{%
  \begin{minipage}[b][\textheight]{\textwidth}
    \vfill
    \usebeamertemplate*{title}
    \ifx\insertsubtitle\@empty\else\usebeamertemplate*{subtitle}\fi
    \usebeamertemplate*{title separator}
    \usebeamertemplate*{author}
    \usebeamertemplate*{institute}
    \vfill
  \end{minipage}}
\makeatother
\setbeamersize{text margin left=1.2em,text margin right=1.2em}

\setbeamertemplate{frame footer}{\insertshorttitle}
\setbeamertemplate{footline}{%
  \begin{beamercolorbox}[wd=\paperwidth,ht=2.6ex,dp=1.2ex,leftskip=1.2em,rightskip=1.2em]{footline}%
    \usebeamerfont{page number in head/foot}%
    \usebeamertemplate*{frame footer}\hfill\usebeamertemplate*{frame numbering}%
  \end{beamercolorbox}}
\setbeamerfont{page number in head/foot}{size=\tiny}
\setbeamercolor{footline}{fg=muted}

% The practical point of a slide, for the baker.
\newcommand{\takeaway}[1]{%
  \par\medskip
  \colorbox{crumb}{\parbox{\dimexpr\linewidth-2\fboxsep}{\small\textbf{\color{crust}In the kitchen:} #1}}}

% Small source line, placed at the bottom of a frame.
\newcommand{\source}[1]{%
  \vfill{\color{muted}\tiny Source: #1}}

% Status flags, mirroring the notes' conventions. Filled tags, so they read on the
% dark title bar as well as on the slide body.
\newcommand{\flagtag}[2]{\mbox{\setlength{\fboxsep}{1.2pt}\colorbox{#1}{\color{white}\scriptsize\bfseries\,#2\,}}}
\newcommand{\hypothesis}{\flagtag{flag}{hypothesis}}
\newcommand{\contested}{\flagtag{flag}{contested}}
\newcommand{\ownresult}{\flagtag{own}{own result}}
\newcommand{\untested}{\flagtag{flag}{not yet tested}}

% Source lists: one source per bullet, author and year in bold, optional group
% headings. Use in a [allowframebreaks] frame; long lists continue on "(cont.)" slides.
\newenvironment{sourcelist}{%
  \setbeamerfont{itemize/enumerate body}{size=\scriptsize}%
  \setbeamerfont{itemize item}{size=\scriptsize}%
  \scriptsize\begin{itemize}\setlength{\itemsep}{2pt}}{\end{itemize}}
\newcommand{\src}[2]{\item\textbf{#1}\enspace #2}
% \nobreak keeps a group heading on the same slide as its first source.
\newcommand{\srcgroup}[1]{\par\smallskip{\small\bfseries\color{crust}#1}\par\nobreak\vspace{-2pt}\nobreak}
\setbeamertemplate{frametitle continuation}[from second][(cont.)]

% Flowchart nodes for TikZ diagrams.
\tikzset{
  box/.style={draw=crust,fill=crumb,rounded corners=2pt,align=center,
    font=\scriptsize,inner sep=4pt,text width=#1},
  box/.default=2.6cm,
  arrow/.style={-{Stealth[length=5pt]},draw=crust,thick},
}

\author{Sourdough Science}
\institute{github.com/zrrbite/sourdough-science}
\date{}


% Local styles for this deck only.
\tikzset{
  nohyph/.style={execute at begin node={\hyphenpenalty=10000
    \exhyphenpenalty=10000}},
  sbox/.style={box=#1,minimum height=2.5em,nohyph},
  sbox/.default=2.6cm,
  hbox/.style={box=#1,fill=crust!18,font=\scriptsize\bfseries,nohyph},
  hbox/.default=2.6cm,
  darrow/.style={arrow,dashed},
  lbl/.style={font=\tiny\itshape,text=muted,align=center},
  head/.style={font=\scriptsize\bfseries,text=crust,align=center},
}
\newcommand{\muted}[1]{{\color{muted}#1}}
\newcolumntype{Y}{>{\raggedright\arraybackslash}X}
\newcommand{\tablefont}{\footnotesize\renewcommand{\arraystretch}{1.15}}

\title[Application]{Application: From Flavour Chemistry to the Bake}
\subtitle{The hybrid dough, pizza, bread and oven spring:\\
  Part~III of the notes}

\begin{document}

\begin{frame}[plain]\titlepage\end{frame}

% ---------------------------------------------------------------
\begin{frame}{What this deck covers}
  Practical questions, answered from Part~III of the notes:
  \medskip
  \begin{columns}[T,onlytextwidth]
    \begin{column}{0.485\textwidth}
      \begin{block}{The hybrid dough}\small
        Why leaven with commercial yeast and add a little levain?
        Where do the fruity esters come from, and what does the oven do
        to them?
      \end{block}
      \begin{block}{Bread}\small
        How do you get the benefits of proteolysis without a long
        ferment? Which lever matters most?
      \end{block}
    \end{column}
    \begin{column}{0.485\textwidth}
      \begin{block}{Pizza}\small
        Why is a longer ferment not always better? Why does pizza burn,
        and how do you run a 12" Ooni?
      \end{block}
      \begin{block}{Handling and oven spring}\small
        What do you do with dough straight from the fridge? What,
        besides gas, decides how much it rises?
      \end{block}
    \end{column}
  \end{columns}
  \source{chapters 12--20 of the notes, and the note \emph{Accelerating
    Proteolysis: pH, Heat, and Salt}}
\end{frame}

% ===============================================================
\section{The hybrid and its three flavour systems}

\begin{frame}{Let commercial yeast leaven and the levain flavour}
  \begin{columns}[T,onlytextwidth]
    \begin{column}{0.52\textwidth}
      \begin{itemize}\setlength{\itemsep}{0.35em}
        \item In 100\% sourdough the wild yeast does all the leavening
        \item That needs a long ferment, so LAB proteases work on the
          gluten for longer
        \item Result: fragile, tear-prone dough that is hard to time
        \item Hybrid: commercial \textit{S.~cerevisiae} makes the gas;
          a small levain brings acid and flavour
      \end{itemize}
    \end{column}
    \begin{column}{0.46\textwidth}
      \centering
      \begin{tikzpicture}
        \node[sbox=2.5cm] (y) at (0,0) {Commercial yeast};
        \node[sbox=2.5cm] (l) at (3.3,0) {Small levain};
        \node[sbox=2.5cm,below=0.55cm of y] (g)
          {Reliable gas,\\predictable schedule};
        \node[sbox=2.5cm,below=0.55cm of l] (f) {Organic acids,\\flavour};
        \node[hbox=5.6cm,below=0.55cm of $(g.south)!0.5!(f.south)$] (h)
          {Shorter ferment, less protease exposure:\\gluten stays intact for
          hand stretching};
        \draw[arrow] (y) -- (g);
        \draw[arrow] (l) -- (f);
        \draw[arrow] (g.south) -- (g.south |- h.north);
        \draw[arrow] (f.south) -- (f.south |- h.north);
      \end{tikzpicture}
    \end{column}
  \end{columns}
  \takeaway{treat leavening and flavour as two separate jobs and tune each
    on its own. The levain works even at a small percentage.}
  \source{ch.~12}
\end{frame}

% ---------------------------------------------------------------
\begin{frame}{Esters are made by enzymes, not by acid meeting alcohol}
  \small
  \begin{itemize}\setlength{\itemsep}{0.1em}
    \item Acid + alcohol → ester + water is the minor route in dough: very slow
      without a catalyst, and the excess water pushes it back
      \muted{(general chemistry; no measurement in dough found)}
    \item Yeast enzymes make their acyl donor inside the cell: the levain's
      acids are not the feedstock, the alcohols are
  \end{itemize}
  \smallskip
  \begingroup\tablefont
  \begin{tabularx}{\linewidth}{@{}l l Y@{}}
    \toprule
    \textbf{Yeast genes} & \textbf{Join} & \textbf{To give} \\
    \midrule
    \textit{ATF1}, \textit{ATF2} & acetyl-CoA + an alcohol &
      ethyl acetate (fruity, solvent-like), isoamyl acetate (banana),
      \mbox{2-phenylethyl} acetate (rose, honey) \\
    \textit{EHT1}, \textit{EEB1} & fatty acyl-CoA + ethanol &
      ethyl hexanoate (apple), ethyl octanoate \\
    \bottomrule
  \end{tabularx}\endgroup
  \takeaway{any dough with active yeast makes these esters, levain or not.
    More fusel alcohol tends to mean more acetate esters; the genes vary by
    strain.}
  \source{Verstrepen et al.\ (2003); Saerens et al.\ (2010); Pradal et al.\
    (2025); ch.~13}
\end{frame}

% ---------------------------------------------------------------
\begin{frame}{Some sourdough LAB make esters too, with the yeast's help}
  \begin{columns}[T,onlytextwidth]
    \begin{column}{0.6\textwidth}
      \small
      \begin{itemize}\setlength{\itemsep}{0.25em}
        \item In a medium built like wheat sourdough, strains of several LAB
          species made ethyl acetate and ethyl lactate
        \item Output was limited by the supply of precursors
        \item Esterase genes vary by species; the \textit{Fructilactobacillus
          sanfranciscensis} strains tested made neither
        \item Ethyl lactate is lactic acid + ethanol, so this route plausibly
          needs both partners
      \end{itemize}
    \end{column}
    \begin{column}{0.38\textwidth}
      \centering
      \begin{tikzpicture}
        \node[sbox=1.9cm] (lab) at (0,0) {LAB:\\lactic acid};
        \node[sbox=1.9cm] (yst) at (2.55,0) {Yeast:\\ethanol};
        \node[hbox=3.4cm,below=0.7cm of $(lab.south)!0.5!(yst.south)$] (el)
          {Ethyl lactate (creamy)\\and ethyl acetate};
        \draw[arrow] (lab.south) -- (lab.south |- el.north);
        \draw[arrow] (yst.south) -- (yst.south |- el.north);
        \node[lbl,below=0.1cm of el] {LAB enzyme, during fermentation:\\the
          one route close to ``acid + alcohol''};
      \end{tikzpicture}
    \end{column}
  \end{columns}
  \takeaway{the levain may add ethyl lactate and more ethyl acetate, but not
    every levain carries LAB that make them.}
  \source{Pradal et al.\ (2025); ch.~13}
\end{frame}

% ---------------------------------------------------------------
\begin{frame}{The oven keeps or loses esters; it does not make them}
  \small
  \begin{itemize}\setlength{\itemsep}{0.2em}
    \item Esters and ethanol are volatile: the crumb stays near 98°C while
      water leaves through a crust past 150°C, carrying them with it
    \item Crust aroma is Maillard and Strecker chemistry: pyrazines,
      furanones, Strecker aldehydes, melanoidins
    \item Pumpernickel: aroma formed mostly in baking (malty 3-methylbutanal,
      2-methylbutanal, potato-like methional); ester levels were set by the
      fermentation \muted{(hours at low heat: extreme, but the same way)}
    \item No source was found for esters forming in the crust
  \end{itemize}
  \takeaway{the fruity character is set while the yeast works. A shorter
    ferment may keep more of it, since CO₂ leaving the dough during folds
    and degassing plausibly carries esters away. \hypothesis}
  \source{Majcher et al.\ (2020); ch.~13}
\end{frame}

% ---------------------------------------------------------------
\begin{frame}{A long cold ferment's complexity is Maillard, not esters}
  \small
  \begin{itemize}\setlength{\itemsep}{0.3em}
    \item 48--72\,h at 3--5°C: flour enzymes keep working, at reduced but
      significant rates
    \item Amylases turn damaged starch into maltose and glucose: roughly
      10--15\% of room-temperature activity, but 48\,h adds up
    \item Flour proteases cut gluten and other proteins into free amino
      acids and short peptides
    \item In the oven these react: melanoidins (crust colour), pyrazines
      (roasted, nutty), furanones (caramel-like), Strecker aldehydes (malty,
      bready, chocolatey)
    \item Yeast's own acids are millimolar, LAB's tens of millimolar: far
      less sour. Its esters give only a low baseline
  \end{itemize}
  \takeaway{the Maillard gain comes from time and enzyme breakdown, so a
    long retard pays off whichever organism does the leavening.}
  \source{ch.~14}
\end{frame}

% ---------------------------------------------------------------
\begin{frame}{The hybrid stacks three flavour systems in one dough}
  \centering
  \begin{tikzpicture}
    \foreach \x/\n/\t in {-5/h1/{1 · Maillard chemistry},
        0/h2/{2 · Ester chemistry}, 5/h3/{3 · Direct LAB metabolites}}
      \node[head] (\n) at (\x,0) {\t};
    \node[sbox=4.1cm,below=0.12cm of h1] (a)
      {Free amino acids + reducing sugars,\\built up in the cold retard};
    \node[sbox=4.1cm,below=0.12cm of h2] (c)
      {Yeast: alcohols + acetyl-CoA\\LAB: lactic acid + ethanol};
    \node[sbox=4.1cm,below=0.12cm of h3] (e)
      {LAB metabolism\\in the levain};
    \node[sbox=4.1cm,below=0.45cm of a] (b)
      {Melanoidins · pyrazines\\furanones · Strecker aldehydes};
    \node[sbox=4.1cm,below=0.45cm of c] (d)
      {Acetate esters · ethyl esters\\ethyl lactate};
    \node[sbox=4.1cm,below=0.45cm of e] (f)
      {Lactic acid · acetic acid\\diacetyl · acetaldehyde};
    \draw[arrow] (a) -- node[lbl,right]{oven heat} (b);
    \draw[arrow] (c) -- node[lbl,right]{enzymes, in fermentation} (d);
    \draw[arrow] (e) -- (f);
    \node[hbox=4.1cm,below=0.35cm of d] (g) {Complete hybrid flavour profile};
    \draw[arrow] (b.south) |- (g.west);
    \draw[arrow] (d) -- (g);
    \draw[arrow] (f.south) |- (g.east);
  \end{tikzpicture}
  \par\raggedright
  \takeaway{commercial yeast alone gives system 1 and the yeast half of 2.
    Pure sourdough gives all three, but with a smaller yeast population, a longer ferment
    that may lose more volatiles, and structural and timing penalties. The
    hybrid keeps all three, with structure and schedule intact.}
  \source{ch.~14}
\end{frame}

% ===============================================================
\section{Pizza: the dough, the long ferment and the Ooni}

\begin{frame}{The oven decides which preferment suits a pizza}
  \centering
  \begin{tikzpicture}
    \node[hbox=1.5cm] (q) at (0,0) {Which oven?};
    \node[sbox=2.9cm] (deck) at (3.35,1.2)
      {Neapolitan deck oven\\450°C+, 60--90\,s bake};
    \node[sbox=2.9cm] (home) at (3.35,-1.1)
      {Home oven\\250--300°C, longer bake};
    \node[sbox=3.1cm] (char) at (6.95,1.2)
      {Char dominates;\\subtle flavour\\differences compress};
    \node[sbox=3.1cm] (acid) at (6.95,-1.1)
      {Crust has time: sourdough can be the \emph{more} flavourful choice,
      if you manage the acid};
    \node[sbox=3.9cm] (pool) at (11.15,1.2)
      {Poolish or direct commercial yeast, 24--72\,h cold: maximum
      extensibility, no acid interference};
    \node[sbox=3.9cm] (young) at (11.15,-0.4)
      {Young liquid levain at peak:\\cornicione puff, less acid};
    \node[sbox=3.9cm] (stiff) at (11.15,-1.8)
      {Stiff levain: slower ferment,\\wider timing window};
    \draw[arrow] (q) |- (deck);
    \draw[arrow] (q) |- (home);
    \draw[arrow] (deck) -- (char);
    \draw[arrow] (home) -- (acid);
    \draw[arrow] (char) -- (pool);
    \draw[arrow] (acid.east) -- (young.west);
    \draw[arrow] (acid.east) -- (stiff.west);
  \end{tikzpicture}
  \par\raggedright
  \takeaway{acid-modified gluten is more extensible but less elastic: less
    snap-back, less resilience if overworked. Reduce hydration slightly or
    raise flour protein.}
  \source{ch.~15}
\end{frame}

% ---------------------------------------------------------------
\begin{frame}{Longer is not automatically better for pizza}
  \centering
  \begin{tikzpicture}
    \node[hbox=5cm] (top) at (0,0) {Extended fermentation:\\48--72\,h and beyond};
    \foreach \x/\n/\c/\e in {
        -6/1/{Proteolysis:\\peptide bonds cut}/{Stretches willingly,\\then tears},
        -3/2/{LAB active at 3--5°C,\\yeast is not}/{48\,h: perfect;\\56\,h: over-proofed},
        0/3/{Acid load\\keeps climbing}/{Slack; loses\\elastic recovery},
        3/4/{CO₂ from degassing carries volatiles out \hypothesis}/{Estery top notes
          fade, acidity rises},
        6/5/{Amylase turns more\\starch to sugar}/{Less starch to crisp:\\gummy crust}} {
      \node[sbox=2.55cm,anchor=north] (c\n) at (\x,-0.8) {\c};
      \node[sbox=2.55cm,anchor=north] (e\n) at (\x,-2.35) {\e};
      \draw[arrow] (c\n) -- (e\n);
      \draw[arrow] (top.south) -- (c\n.north);
    }
    \node[hbox=7cm] (end) at (0,-3.85)
      {Diminishing returns, then losses};
    \foreach \n in {1,...,5} \draw[arrow] (e\n.south) -- (end.north);
  \end{tikzpicture}
  \par\raggedright
  \takeaway{many excellent pizzas use 48--72\,h cold doughs, but longer is
    not better by default. A shorter, controlled ferment with a hybrid's
    flavour can be more reliable, and arguably more complex.}
  \source{ch.~16; CO₂ stripping of volatiles is a hypothesis (ch.~13)}
\end{frame}

% ---------------------------------------------------------------
\begin{frame}{Over-fermented pizza dough fails suddenly, not gradually}
  \small
  \begin{itemize}\setlength{\itemsep}{0.3em}
    \item LAB proteases, and to a lesser extent flour proteases, keep
      cutting the gluten network
    \item A loaf can use that; a pizza must survive being stretched by hand
      into a thin, continuous sheet
    \item Degraded dough stretches willingly until it suddenly doesn't: thin
      spots and holes, worst in the thin centre
    \item Acid: moderate acidity helps stretching, too much turns it slack;
      acetic acid feels tight, then suddenly gives way
    \item LAB keep working at 3--5°C while yeast nearly stops: perfect at
      48\,h, over-proofed at 56\,h
  \end{itemize}
  \takeaway{with sourdough the usable window is narrow and hard to predict.
    Low-dose commercial yeast in a 48--72\,h cold ferment can give a window
    several hours wide.}
  \source{ch.~16}
\end{frame}

% ---------------------------------------------------------------
\begin{frame}{Residual sugar, not hydration, is why pizza burns}
  \begingroup\tablefont
  \begin{tabularx}{\linewidth}{@{}l Y Y@{}}
    \toprule
    \textbf{Ferment} & \textbf{Sugar at launch} & \textbf{At 450°C+} \\
    \midrule
    Short & yeast has not eaten the sugar pool & uniform carbonisation in
      seconds \\
    24--72\,h cold retard & yeast has eaten most simple sugars & controlled
      leopard spots \\
    Too long & amylase outruns the yeast: sugars re-accumulate &
      burns \emph{and} handles poorly \\
    \bottomrule
  \end{tabularx}\endgroup
  \small
  \begin{itemize}\setlength{\itemsep}{0.1em}
    \item Flour has about 1.5--2\% free sugars, plus damaged starch that
      amylases turn into maltose and glucose
    \item Water pins the surface near 100°C until it dries: a large buffer
      at 250--300°C, mere seconds at 450°C+
  \end{itemize}
  \takeaway{fix burning through the ferment, not the water: more hydration
    buys seconds but cannot rescue a sugar-rich dough.}
  \source{ch.~18}
\end{frame}

% ---------------------------------------------------------------
\begin{frame}{In a 12" Ooni, two ball sizes both work}
  \begingroup\tablefont\renewcommand{\arraystretch}{1.0}
  \begin{tabularx}{\linewidth}{@{}l Y Y@{}}
    \toprule
    & \textbf{Larger ball: forgiving} & \textbf{Smaller ball: crisp} \\
    \midrule
    Ball → pie & 270--280\,g → 10--11" & 220--230\,g → 11--12" \\
    Crust & thicker, more thermal mass & thinner, cooks through fast \\
    At 40--50\% flame & 2--3\,min bake & 90\,s to 2\,min bake \\
    Margin for error & wide & small: done to carbonised in seconds \\
    Maillard aroma & more time & less time \\
    Eats like & bread-tender, crisp outside; Roman/NY & cracker-crisp, no
      raw centre \\
    \bottomrule
  \end{tabularx}\endgroup
  \smallskip
  \small
  \begin{itemize}
    \item Hydration 63--67\%: above about 68--70\% the dough sticks on launch
      and sags, and if it folds on the stone there is no room to recover
  \end{itemize}
  \takeaway{the choice is stylistic. Both want a 48\,h cold retard (sugar
    depletion) and a moderate 40--50\% of max flame.}
  \source{ch.~18}
\end{frame}

% ---------------------------------------------------------------
\begin{frame}{Twice the crust thickness, four times the time to cook through}
  \begin{columns}[T,onlytextwidth]
    \begin{column}{0.67\textwidth}
      \small
      \begin{itemize}\setlength{\itemsep}{0.1em}
        \item Heat comes from stone and flame and diffuses slowly inward
        \item Neapolitan crust (\textasciitilde3\,mm): gelatinised (60--80°C)
          right through, well within 60\,s
        \item 275\,g ball at 10--11": \textasciitilde6--8\,mm, the path
          roughly doubles
        \item Time scales with distance squared (Fick's law): about 4×
        \item Full flame chars the surface first: burnt outside, gummy or raw
          inside
      \end{itemize}
    \end{column}
    \begin{column}{0.31\textwidth}
      \centering
      \begin{tikzpicture}
        % thin crust
        \fill[crumb,draw=crust] (0,2.2) rectangle (1.8,2.5);
        \foreach \x in {0.3,0.9,1.5} {
          \draw[arrow,flag] (\x,2.95) -- (\x,2.55);
          \draw[arrow,flag] (\x,1.75) -- (\x,2.15);
        }
        \node[font=\scriptsize,anchor=west,align=left] at (1.9,2.35)
          {\textasciitilde3\,mm\\\muted{through in}\\\muted{seconds}};
        % thick crust
        \fill[crumb,draw=crust] (0,0) rectangle (1.8,0.75);
        \foreach \x in {0.3,0.9,1.5} {
          \draw[arrow,flag] (\x,1.2) -- (\x,0.8);
          \draw[arrow,flag] (\x,-0.45) -- (\x,-0.05);
        }
        \node[font=\scriptsize,anchor=west,align=left] at (1.9,0.375)
          {\textasciitilde6--8\,mm\\\muted{2× the path,}\\\muted{\textasciitilde4×
          the time}};
        \node[lbl,anchor=north] at (0.9,-0.55) {flame above,\\stone below};
      \end{tikzpicture}
    \end{column}
  \end{columns}
  \takeaway{trade surface heat for diffusion time: lower flame, longer bake,
    so the centre cooks before the outside chars.}
  \source{Sadot et al.\ (2017); ch.~18}
\end{frame}

% ---------------------------------------------------------------
\begin{frame}{Ooni 12": preheat hard, bake moderate, boost at the end}
  \begingroup\tablefont
  \begin{tabularx}{\linewidth}{@{}l l Y@{}}
    \toprule
    \textbf{Step} & \textbf{Flame} & \textbf{Target} \\
    \midrule
    Preheat & full, about 15\,min & stone heat-soaked; IR reading 380--420°C.
      Longer risks an overheated thin stone \\
    Pre-launch & 40--50\% of max & stone 350--380°C; check 30--60\,s after
      the flame settles \\
    Bake & 40--50\% & 2--3\,min, turning every 20--30\,s \\
    Char boost & briefly up & final 15--20\,s, optional: leopard spots on
      the cornicione \\
    Recovery & full & stone drops 50--100°C per pizza; 3--5\,min back to
      target, then 40--50\% \\
    \bottomrule
  \end{tabularx}\endgroup
  \takeaway{this suits the hybrid. The retard leaves enough sugar to brown
    but not burn; 2--3\,min gives more Maillard and Strecker time than a
    60\,s bake; the esters were set in fermentation. A cold stone gives a
    pale, soggy bottom.}
  \source{ch.~18}
\end{frame}

% ===============================================================
\section{Bread: controlled proteolysis}

\begin{frame}{For bread, proteolysis is a tool rather than a failure}
  \begingroup\tablefont
  \begin{tabularx}{\linewidth}{@{}l Y Y@{}}
    \toprule
    & \textbf{Pizza} & \textbf{Bread} \\
    \midrule
    Failure mode & tears when stretched to a 3--6\,mm sheet & loses gas
      retention: flat, dense loaf \\
    Threshold & near & much further out \\
    Tolerable levain & 10--15\% & 20--30\% \\
    \bottomrule
  \end{tabularx}\endgroup
  \smallskip
  \small
  \begin{itemize}\setlength{\itemsep}{0.1em}
    \item \textbf{Extensibility:} easier shaping, more oven spring before
      the network resists
    \item \textbf{Softer crumb:} partly depolymerised glutenin bites less
      chewy
    \item \textbf{Free amino acids:} Maillard substrate, fully expressed in
      a 40--50\,min bake, near-irrelevant in a 2--3\,min pizza bake
  \end{itemize}
  \takeaway{a loaf takes far more proteolysis than a pizza before it fails:
    use it on purpose.}
  \source{Bombara et al.\ (1997); ch.~17; note \emph{Accelerating
    Proteolysis}}
\end{frame}

% ---------------------------------------------------------------
\begin{frame}{Acidifying the levain barely acidifies the dough}
  \begin{columns}[T,onlytextwidth]
    \begin{column}{0.56\textwidth}
      \small
      \begin{itemize}\setlength{\itemsep}{0.3em}
        \item Cereal aspartic proteases peak around pH 3.5--4.0
        \item Buffering is molar: pH 6.0 → 4.0 needs a fixed quantity of
          acid, however fast it is made
        \item Hot LAB (40--45°C) acidify perhaps 2--4× faster: 8\,h might
          become 2--3\,h, not 30\,min
        \item Higher-ash flour buffers harder: its minerals resist the drop
          that would activate its proteases
      \end{itemize}
    \end{column}
    \begin{column}{0.42\textwidth}
      \centering
      \begin{tikzpicture}[node distance=0.42cm]
        \node[sbox=3.6cm] (l) {Hot levain reaches pH 3.8};
        \node[sbox=3.6cm,below=of l] (m)
          {Mixed in at 10--15\% of total flour};
        \node[sbox=3.6cm,below=of m] (d)
          {Bulk flour buffers it back: dough pH 5.2--5.5};
        \node[hbox=3.6cm,below=of d] (r)
          {Bulk flour's own proteases stay largely dormant};
        \draw[arrow] (l) -- (m);
        \draw[arrow] (m) -- (d);
        \draw[arrow] (d) -- (r);
      \end{tikzpicture}
    \end{column}
  \end{columns}
  \takeaway{getting the levain acidic is easy; getting the final dough to
    pH 3.5--4.0 quickly is not, and heat alone will not do it.}
  \source{Thiele et al.\ (2004); ch.~17}
\end{frame}

% ---------------------------------------------------------------
\begin{frame}{In a hot bulk, heat does the work that pH cannot}
  \begin{columns}[T,onlytextwidth]
    \begin{column}{0.58\textwidth}
      \small
      \begin{itemize}\setlength{\itemsep}{0.3em}
        \item Enzyme rates roughly double for every 10°C rise
          (Q₁₀\,≈\,2), up to denaturation
        \item Wheat proteases stay stable to about 50°C
        \item Further from the pH optimum but nearer the thermal optimum:
          the two effects partly cancel
        \item Not a dramatic acidification: ordinary proteases running hot
      \end{itemize}
    \end{column}
    \begin{column}{0.4\textwidth}
      \centering
      \begin{tikzpicture}
        \node[sbox=2.3cm] (a) at (0,0) {pH 4.5\\25°C};
        \node[sbox=2.3cm] (b) at (3.3,0) {pH 5.5\\40°C};
        \node[font=\Large\bfseries,text=crust] at (1.65,0) {≈};
        \node[lbl] at (1.65,-0.9) {comparable total proteolytic\\throughput};
      \end{tikzpicture}
    \end{column}
  \end{columns}
  \bigskip
  {\footnotesize\muted{Caveat: Q₁₀\,≈\,2 is a general biochemical
    approximation, not a bread measurement. Use it as an order-of-magnitude
    guide.}}
  \takeaway{don't chase a low pH; use the temperature axis of the enzyme
    curve instead.}
  \source{Bloksma \& Bushuk (1988); ch.~17}
\end{frame}

% ---------------------------------------------------------------
\begin{frame}{Salt is the biggest brake on proteolysis}
  \small
  \begin{itemize}\setlength{\itemsep}{0.35em}
    \item Salt suppresses cereal proteases two ways: higher ionic strength
      (enzyme meets substrate less well) and lower water activity (less
      hydrolysis)
    \item It works regardless of pH, and the effect is substantial at normal
      bread levels of 1.8--2.2\% of flour
    \item An unsalted dough lets its proteases run considerably faster
    \item Calvel's autolyse is usually explained as hydration. That is
      correct but incomplete: part of its softening is enzymatic
  \end{itemize}
  \takeaway{hold the salt back until the final mix. It is a bigger
    proteolytic lever than any realistic pH change, costs nothing, and is
    reversed by salting.}
  \source{Beck et al.\ (2012); Calvel et al.\ (2001); ch.~17}
\end{frame}

% ---------------------------------------------------------------
\begin{frame}{A warm autolyse gives proteolysis in an hour, no fermentation}
  \centering
  \begin{tikzpicture}
    \node[sbox=2.5cm] (a) at (0,0) {Flour + all water at 40--42°C;
      no salt, levain or yeast};
    \node[sbox=2.3cm] (b) at (3.05,0) {Hold 60--90\,min at 38--42°C};
    \node[hbox=2.3cm] (c) at (6.1,0) {Assess: soft, stretches thin?};
    \node[sbox=2.3cm] (d) at (9.15,0) {Final mix: add levain and salt};
    \node[sbox=2.3cm] (e) at (12.2,0) {Salt closes the window; bulk and
      shape as normal};
    \node[sbox=2.9cm] (t) at (6.1,1.2) {Too tight: extend 30\,min, re-assess};
    \node[sbox=2.9cm] (s) at (6.1,-1.2) {Too slack: less time or heat next bake};
    \draw[arrow] (a) -- (b);
    \draw[arrow] (b) -- (c);
    \draw[arrow] (c) -- node[lbl,above]{yes} (d);
    \draw[arrow] (d) -- (e);
    \draw[arrow] ([xshift=-4mm]c.north) -- ([xshift=-4mm]t.south);
    \draw[arrow] ([xshift=4mm]t.south) -- ([xshift=4mm]c.north);
    \draw[arrow] (c.south) -- (s.north);
  \end{tikzpicture}
  \par\smallskip\raggedright\small
  \begin{itemize}\setlength{\itemsep}{0em}
    \item Proteases at pH \textasciitilde6.0 but 2--4× faster (Q₁₀), no salt:
      no acid, buffering or microbes to manage
    \item Stay below about 50°C: starch gelatinises from 55--60°C, making a
      scald (\textit{Brühstück}) that denatures the proteases
    \item Thermal lag: a large mass in a plastic tub takes 30--45\,min to
      reach the setpoint at its core
  \end{itemize}
  \source{Eliasson (1980); ch.~17. \untested{} against a room-temperature
    control}
\end{frame}

% ---------------------------------------------------------------
\begin{frame}{Rank the levers: salt first, pH last}
  \begingroup\tablefont
  \begin{tabularx}{\linewidth}{@{}c l Y@{}}
    \toprule
    & \textbf{Lever} & \textbf{Why} \\
    \midrule
    1 & Remove salt from the autolyse & the largest single effect; free;
      reversed by salting at the final mix \\
    2 & Raise autolyse or bulk to 38--42°C & doubles to quadruples enzyme
      rate via Q₁₀, no pH change needed \\
    3 & Increase the levain percentage & moves acid and LAB protease load
      together; bread tolerates 20--30\%, but the loaf gets more acidic \\
    4 & Extend time & always works, but it is what we set out to avoid \\
    5 & Manipulate pH directly & least efficient via fermentation: buffering
      and dilution \\
    \bottomrule
  \end{tabularx}\endgroup
  \takeaway{with fermentation and proteolysis decoupled, the ferment can be
    as long or as short as the flavour goal needs.}
  \source{ch.~17}
\end{frame}

% ---------------------------------------------------------------
\begin{frame}{Direct acidification: a shortcut that tightens before it
    softens}
  \small
  \begin{itemize}\setlength{\itemsep}{0.3em}
    \item A little food-grade lactic acid reaches the target pH at once, with
      no wait for LAB
    \item With a 40°C hold: rapid proteolysis with no extended fermentation
    \item But added acid first firms the dough and cuts extensibility
      (charge-state and hydration changes)
    \item Held warm, enzymatic softening catches up and overtakes; expect a
      phase where the dough seems to go the wrong way
    \item Chemical and biological acidification do not give the same
      rheology: the literature here is contested \contested
  \end{itemize}
  \takeaway{a specialist tool. For most bread, the no-salt warm autolyse
    gives enough proteolysis with fewer variables and no risk of a dough
    that cannot hold gas.}
  \source{Clarke et al.\ (2004); Wehrle et al.\ (1997); ch.~17; ch.~25}
\end{frame}

% ---------------------------------------------------------------
\begin{frame}{The promised flavour dividend is the least-evidenced claim}
  \small
  \begin{itemize}\setlength{\itemsep}{0.35em}
    \item The claim: the warm autolyse enlarges the free amino acid pool
    \item In a 40--50\,min bake that pool feeds Strecker degradation
    \item Expected result: malty, nutty and roasted crust notes from
      Strecker aldehydes and pyrazines
    \item The payoff a short pizza bake cannot access
  \end{itemize}
  \medskip
  \colorbox{crumb}{\parbox{\dimexpr\linewidth-2\fboxsep}{\small
    \untested{} The repo calls this the least directly evidenced claim in
    the document. It wants a sensory comparison: warm-autolyse loaf against
    a control on the same formula.}}
  \source{Strecker chemistry: Yang et al.\ (2022); claim: ch.~17 and note
    \emph{Accelerating Proteolysis}; status: README}
\end{frame}

% ===============================================================
\section{Handling and oven spring}

\begin{frame}{After the retard, pull the dough; don't press it}
  \begin{columns}[T,onlytextwidth]
    \begin{column}{0.44\textwidth}
      \centering
      \begin{tikzpicture}[node distance=0.42cm]
        \node[sbox=4.6cm] (a) {Retard at 3--5°C: more CO₂ dissolves (Henry's
          law)};
        \node[sbox=4.6cm,below=of a] (b) {Tempering warms the dough: the CO₂
          is supersaturated};
        \node[sbox=4.6cm,below=of b] (c) {It needs nucleation sites; a still
          ball has few fresh surfaces};
        \node[hbox=4.6cm,below=of c] (d) {A gentle pull makes micro-voids:
          many small new bubbles};
        \draw[arrow] (a) -- (b);
        \draw[arrow] (b) -- (c);
        \draw[arrow] (c) -- (d);
      \end{tikzpicture}
    \end{column}
    \begin{column}{0.54\textwidth}
      \small
      \begin{itemize}\setlength{\itemsep}{0.3em}
        \item Like tapping a carbonated drink: disturbance triggers
          nucleation, not more yeast activity
        \item The pull also re-tensions the gluten, which holds the bubbles
          and lifts vertically
        \item Pressing flat expels gas; sluggish post-retard yeast must
          re-inflate from scratch
        \item Pull \emph{adds} gas; press \emph{removes} it
      \end{itemize}
    \end{column}
  \end{columns}
  \takeaway{lift an edge, stretch without tearing, fold over the centre;
    rotate 90° and repeat 2--3 times. Re-round loosely, don't degas. Expect
    fast rise in 20--30\,min; tempering stays 1.5--2\,h.}
  \source{ch.~19; CO₂ partitioning: Chiotellis \& Campbell (2003)}
\end{frame}

% ---------------------------------------------------------------
\begin{frame}{Same dough, three handlings, three different rises
    \ownresult}
  \begingroup\tablefont
  \begin{tabularx}{\linewidth}{@{}c l Y l@{}}
    \toprule
    & \textbf{Retarded in} & \textbf{Then} & \textbf{Result} \\
    \midrule
    1 & wide, flat container & turned out, cut into rolls & good rise \\
    2 & bowl & turned out, cut into rolls & less rise \\
    3 & wide (dough of 1) & flattened like a pizza centre, slapped, baked
      flat & great rise, open crumb \\
    \bottomrule
  \end{tabularx}\endgroup
  \smallskip
  \small
  \begin{itemize}\setlength{\itemsep}{0.15em}
    \item Same formula, retard and yeast activity: only the mechanical state
      of the gluten changed
    \item Wide container: dough spreads under gravity, strands align
      horizontally (laminar), gas finds an easier path upward
    \item Bowl: a tangled, random mesh with a compressed bottom, resisting
      expansion in every direction
  \end{itemize}
  \takeaway{retard in a wide, flat container rather than a bowl.}
  \source{ch.~20; the author's own experiments}
\end{frame}

% ---------------------------------------------------------------
\begin{frame}{Steam delamination: gluten alignment beats gas volume}
  \begin{columns}[T,onlytextwidth]
    \begin{column}{0.53\textwidth}
      \small
      \begin{itemize}\setlength{\itemsep}{0.25em}
        \item A thin disc reaches steam temperature right through at once;
          a roll heats as a slow wave
        \item Water turns to steam, about 1,600-fold in volume
        \item Steam follows the weak interfaces between aligned gluten
          planes and pushes them apart in one event: pita physics
        \item The holes are gaps between separated layers, not coalesced
          fermentation bubbles
      \end{itemize}
    \end{column}
    \begin{column}{0.45\textwidth}
      \centering
      \begin{tikzpicture}[node distance=0.25cm]
        \node[sbox=3.4cm] (a) {1 · Gluten alignment\\(dominant)};
        \node[sbox=3.4cm,below=of a] (b) {2 · Heating uniformity};
        \node[sbox=3.4cm,below=of b] (c) {3 · Network extensibility};
        \node[sbox=3.4cm,below=of c] (d) {4 · Gas volume, nucleation\\
          (secondary)};
        \node[hbox=1.95cm,right=0.4cm of $(b.east)!0.5!(c.east)$] (e)
          {Explosive delamination:\\oven spring};
        \foreach \n in {a,b,c,d} \draw[darrow] (\n.east) -- (e.west);
      \end{tikzpicture}
    \end{column}
  \end{columns}
  \takeaway{a well-aligned, largely degassed dough outperforms a gassy dough
    with random gluten orientation. Thinner pieces spring more per gram.}
  \source{ch.~20; the hierarchy is drawn from the author's own experiments
    \ownresult}
\end{frame}

% ---------------------------------------------------------------
\begin{frame}{Fold to multiply the layers, then rest 15--30 minutes}
  \begin{columns}[T,onlytextwidth]
    \begin{column}{0.485\textwidth}
      \textbf{\color{crust}Fold}\par\small
      \begin{itemize}\setlength{\itemsep}{0.2em}
        \item Letter fold in thirds, rotate 90°, flatten, fold again: 2
          cycles \textasciitilde9 layers, 3 cycles \textasciitilde27
        \item Puff pastry without the fat: rougher open crumb, not flakes
        \item Plain dough blurs at about 20--40 effective layers (puff
          pastry: 700--900)
      \end{itemize}
    \end{column}
    \begin{column}{0.485\textwidth}
      \textbf{\color{crust}Rest}\par\small
      \begin{itemize}\setlength{\itemsep}{0.2em}
        \item Just folded, the gluten is elastic and would contract against
          the steam
        \item Resting relaxes it (thiol--disulfide interchange) while CO₂
          nucleates into the new interfaces
        \item Beyond 30--40\,min the layers merge and the dough spreads
      \end{itemize}
    \end{column}
  \end{columns}
  \takeaway{for rolls or flatbreads: turn out from a wide container, flatten,
    letter-fold twice, portion loosely, rest 15--20\,min until slightly
    puffed, bake at high heat.}
  \source{ch.~20}
\end{frame}

% ---------------------------------------------------------------
\begin{frame}{In pizza, shaping makes the cornicione puff}
  \small
  \begin{itemize}\setlength{\itemsep}{0.35em}
    \item Hand-stretching aligns the gluten laminarly in the rim; oven steam
      then delaminates those layers
    \item A well-aligned rim on slightly under-proofed dough out-puffs a
      poorly handled rim on perfectly proofed dough
    \item Neapolitan slapping is not just shaping the disc: it builds that
      alignment
    \item The pressed-thin centre has too little cross-section to puff; the
      rim has both alignment and thermal mass
  \end{itemize}
  \takeaway{for the hybrid dough in an Ooni: gentle pull, temper 1.5--2\,h,
    then open with deliberate attention to the rim's alignment. Shaping
    matters more than fermentation state for the puff.}
  \source{ch.~20, applying the author's roll and flatbread experiments to
    pizza}
\end{frame}

% ===============================================================
% Wrap-up

\begin{frame}{Takeaways}
  \small
  \begin{itemize}\setlength{\itemsep}{0.25em}
    \item Let commercial yeast leaven and a small levain flavour: three
      flavour systems, intact gluten, a predictable schedule
    \item Esters are made by enzymes while the yeast works; the oven only
      keeps or loses them. Crust aroma is Maillard and Strecker chemistry
    \item For pizza, longer is not automatically better: proteolysis, acid,
      a narrowing window and returning sugar all turn against you
    \item Burning is residual sugar. In a 12" Ooni: 48\,h retard, 63--67\%
      hydration, 40--50\% flame, 2--3\,min on a 350--380°C stone
    \item For bread proteolysis, hold back the salt and warm the autolyse to
      38--42°C (below 50°C) before chasing pH
    \item After the retard, pull, don't press. For oven spring, gluten
      alignment beats gas volume
  \end{itemize}
  \source{chapters 12--20}
\end{frame}

% ---------------------------------------------------------------
\begin{frame}{Open questions}
  \small
  \begin{itemize}\setlength{\itemsep}{0.45em}
    \item Warm autolyse against a room-temperature control on the same
      formula: extensibility by hand, crumb softness after the bake
      \untested
    \item Crust aroma, warm-autolyse loaf vs control: a sensory check of the
      Strecker/pyrazine claim, the least directly evidenced in the document
      \untested
    \item Log the core dough temperature during the warm autolyse to measure
      the thermal lag \untested
    \item Direct lactic acid + 40°C hold as a comparison arm: expect
      tightening before the enzymes soften it \untested
    \item Does CO₂ leaving during folds and degassing strip esters? No
      measurement in dough was found \hypothesis
  \end{itemize}
  \source{note \emph{Accelerating Proteolysis} (open threads); README
    (status); ch.~13; ch.~16}
\end{frame}

% ---------------------------------------------------------------
\begin{frame}[allowframebreaks]{Sources}
  \srcgroup{Esters and fermentation volatiles}
  \begin{sourcelist}
    \src{Verstrepen KJ et al.\ (2003)}{Expression levels of the yeast alcohol acetyltransferase genes \textit{ATF1}, \textit{Lg-ATF1}, and \textit{ATF2} control the formation of a broad range of volatile esters. \textit{Appl Environ Microbiol} 69(9):5228--5237.}
    \src{Saerens SMG, Delvaux FR, Verstrepen KJ, Thevelein JM (2010)}{Production and biological function of volatile esters in \textit{Saccharomyces cerevisiae}. \textit{Microb Biotechnol} 3(2):165--177.}
    \src{Pradal I, Weckx S, De Vuyst L (2025)}{The production of esters by specific sourdough lactic acid bacteria species is limited by the precursor concentrations. \textit{Appl Environ Microbiol} 91(3):e0221624.}
    \src{Majcher MA et al.\ (2020)}{Formation of key aroma compounds during preparation of pumpernickel bread. \textit{J Agric Food Chem} 68(38):10352--10360.}
    \src{Yang Y et al.\ (2022)}{Research progress on the formation mechanism and detection technology of bread flavor. \textit{J Food Sci} 87(8):3330--3351.}
  \end{sourcelist}
  \srcgroup{Gas and heat}
  \begin{sourcelist}
    \src{Chiotellis E, Campbell GM (2003)}{Proving of bread dough I: modelling the evolution of the bubble size distribution. \textit{Food Bioprod Process} 81(3):194--206.}
    \src{Sadot M et al.\ (2017)}{Numerical modelling of convective heat transfer during bread baking. \textit{J Food Eng}.}
  \end{sourcelist}
  \srcgroup{Acidification, proteolysis, salt and autolyse}
  \begin{sourcelist}
    \src{Thiele C, Grassl S, Gänzle MG (2004)}{Gluten hydrolysis and depolymerization during sourdough fermentation. \textit{J Agric Food Chem} 52(5):1307--1314.}
    \src{Clarke CI et al.\ (2004)}{Wheat sourdough fermentation: effects of time and acidification on fundamental rheological properties. \textit{Cereal Chem} 81(3):409--417.}
    \src{Wehrle K, Grau H, Arendt EK (1997)}{Effects of lactic acid, acetic acid, and table salt on fundamental rheological properties of wheat dough. \textit{Cereal Chem} 74(6):739--744.}
    \src{Calvel R, Wirtz RL, MacGuire JJ (2001)}{\textit{The Taste of Bread}. Aspen Publishers.}
    \src{Beck M, Jekle M, Becker T (2012)}{Impact of sodium chloride on wheat flour dough for yeast-leavened products. \textit{J Sci Food Agric} 92(2):299--306.}
    \src{Bombara N, Añón MC, Pilosof AMR (1997)}{Functional properties of protease modified wheat flours. \textit{LWT} 30(5):441--447.}
    \src{Bloksma AH, Bushuk W (1988)}{Rheology and chemistry of dough. In \textit{Wheat: Chemistry and Technology}, 3rd ed., AACC.}
    \src{Eliasson AC (1980)}{Effect of water content on the gelatinization of wheat starch. \textit{Starch/Stärke} 32(8):270--272.}
  \end{sourcelist}
  \medskip
  {\scriptsize\color{muted}Full list: ch.~25 of the notes.\par}
\end{frame}

\end{document}
